\subsection{Software and Firmware Architecture}

\subsubsection{Firmware Architecture}

The embedded program was designed as a control loop in the ESP32-S3 processor. A single
loop configuration was adopted to make the firmware reliable, portable, and easier to test
during multiple field tests. Rather than adopting an RTOS-based multi-threaded program
configuration, the system implements sensing, decision-making, and actuation as different
software modules run sequentially according to a set sequence of operations.

The control loop executes the following functions:
\begin{itemize}
\item Collection of a batch of RSSI measurements from the received data packets.
\item Filtering and averaging of the acquired RSSI samples.
\item Encapsulation of the current pan and tilt angles and RSSI measurement trend.
\item Decision on the next step based on the chosen alignment strategy.
\item Stepper motor and servo command execution.
\item Logging of experiment parameters for later analysis.
\end{itemize}

The above sequence ensures that each alignment algorithm operates on identical feedback
measurements and faces the same delay between successive movements.d the same settling delay between consecutive movements.


\subsubsection{RSSI Processing and State Encoding}

The RSSI signal is highly corrupted by multipath effects, small-scale fading, and
fluctuations across successive packets. This is addressed by averaging a few packets
collected for each orientation before running the decision block. The smoothed RSSI
signal at step $k$ of the control algorithm is obtained as follows:

\begin{equation}
\overline{\text{RSSI}}_k = \frac{1}{N} \sum_{i=1}^{N} \text{RSSI}_{k,i}
\end{equation}
where $N$ is the number of samples collected for the current orientation. The trend
of the RSSI signal over a short period of time can be computed as:
\begin{equation}
\Delta \text{RSSI}_k = \overline{\text{RSSI}}_k - \overline{\text{RSSI}}_{k-1}
\end{equation}
This variable is discretized into three states based on whether the RSSI value
is increasing, constant, or decreasing using a dead-band of size $\epsilon$.

The embedded state is encoded as:
\begin{equation}
s_k = \left( i_{\text{pan}},\; i_{\text{tilt}},\; \text{trend}_k \right)
\end{equation}
where $i_{\text{pan}}$ and $i_{\text{tilt}}$ are discrete orientation indices derived from
fixed angular step sizes.

\subsubsection{Firmware Modules}

The firmware was structured in modules for the exhaustive search, hill climbing, and RL policies to use the same sensing and actuation primitives. This mitigated bias in comparing the algorithms.

\vspace{0.5cm}
\begin{table}[H]
\centering
\caption{Embedded Firmware Modules}
\label{tab:firmware_modules}
\renewcommand{\arraystretch}{1.4}
\begin{tabular}{|
>{\raggedright\arraybackslash}p{3.2cm}|
>{\raggedright\arraybackslash}p{4.2cm}|
>{\raggedright\arraybackslash}p{5.2cm}|
}
\hline
\textbf{Module} & \textbf{Implemented On} & \textbf{Function} \\
\hline
RSSI Acquisition & ESP32 Wi-Fi interface & Captures packet RSSI values and stores a short measurement window for averaging \\
\hline
Motion Control & GPIO and MCPWM peripherals & Generates STEP/DIR pulses for pan motion and PWM pulses for tilt adjustment \\
\hline
State Encoder & Application layer firmware & Maps continuous angles and RSSI trend to discrete state indices \\
\hline
Policy Executor & Application layer firmware & Runs exhaustive scan, hill-climbing, or greedy Q-table inference \\
\hline
Data Logger & UART serial output & Records trial identifier, angles, RSSI, elapsed time, and convergence status \\
\hline
\end{tabular}
\end{table}

\subsubsection{Control Timing and Runtime Parameters}

An equal timing setup was kept throughout all the algorithms to make comparisons
more accurate for the practical testing process. In addition, the firmware waited for
mechanical settling after each movement before taking another RSSI sampling session.
Such a measure was required to prevent possible measurement error due to vibrations
and misalignment.

The selected runtime parameters correspond to the actual findings made using the
prototype. Specifically, the pan range was limited by the visible front hemisphere of
the transmitter, tilt range was determined by mechanical considerations of the
mounting bracket used, and 180~ms delay between movements allowed enough time to
prevent oscillation before sampling.

\vspace{0.5cm}
\begin{table}[H]
\centering
\caption{Key Runtime Parameters Used in Practical Experiments}
\label{tab:runtime_parameters}
\renewcommand{\arraystretch}{1.4}
\begin{tabular}{|
>{\raggedright\arraybackslash}p{3.2cm}|
>{\raggedright\arraybackslash}p{3.2cm}|
>{\raggedright\arraybackslash}p{6.2cm}|
}
\hline
\textbf{Parameter} & \textbf{Value} & \textbf{Remarks} \\
\hline
Pan sweep range & $0^{\circ}$ to $180^{\circ}$ & Limited to the mechanically reachable front hemisphere during trials \\
\hline
Tilt sweep range & $25^{\circ}$ to $75^{\circ}$ & Selected to cover the practical mounting geometry \\
\hline
Coarse pan step & $10^{\circ}$ & Used for exhaustive scan and initial search stages \\
\hline
Coarse tilt step & $5^{\circ}$ & Matches the servo control resolution used in testing \\
\hline
Fine step size & $2^{\circ}$ & Used for local refinement around promising states \\
\hline
Samples per point & 20 & Average computed from 20 consecutive RSSI readings \\
\hline
Sample delay & 10~ms & Delay between individual RSSI samples within a batch \\
\hline
Settling delay after movement & 180~ms & Allows mechanical oscillation to decay before measurement \\
\hline
\end{tabular}
\end{table}

\subsubsection{Experiment Logging Structure}

To aid post-analysis tasks, every practice run was logged using an experiment header that provided information about the test environment and scan settings. The log data structure remained consistent for the entire project duration, making the collected data compatible with spreadsheet and scripting analysis tools.

\vspace{0.5cm}
\begin{table}[H]
\centering
\caption{Practical Experiment Log Header}
\label{tab:experiment_log_header}
\renewcommand{\arraystretch}{1.3}
\begin{tabular}{|
>{\raggedright\arraybackslash}p{3.8cm}|
>{\raggedright\arraybackslash}p{9.4cm}|
}
\hline
\textbf{Field} & \textbf{Description} \\
\hline
Experiment\_ID & Unique identifier of the trial set \\
\hline
Environment & Indoor or outdoor deployment condition \\
\hline
Distance\_TX\_RX\_m & Separation between transmitter and receiver antennas \\
\hline
Obstacle\_Description & LOS, partial blockage, corridor wall, vegetation, and similar descriptors \\
\hline
Pan\_Min/Max, Tilt\_Min/Max & Mechanical search limits used for the run \\
\hline
Pan\_Step, Tilt\_Step, Fine\_Step & Control resolution used by the selected algorithm \\
\hline
Steps\_Per\_Degree & Conversion factor between logical angle and motor command \\
\hline
Samples\_Per\_Point, Sample\_Delay\_ms & RSSI averaging settings for each orientation \\
\hline
Notes & Additional observations recorded during the field run \\
\hline
\end{tabular}
\end{table}
